Dieser Bericht untersucht, wo die im September 2026 veröffentlichte \emph{Graph Zeta Library} (GZL) über ihren ursprünglichen Zweck hinaus neue Ergebnisse liefern kann. Grundlage sind eine Recherche auf arXiv, vor allem zu Arbeiten aus 2025 und 2026, und eigene Rechnungen. Die stärkste Idee ist ein Atlas der Gitter, die Mehrkörper-Potenzgesetzenergien minimieren. Die zugehörigen Pilotrechnungen wurden mit einer von GZL unabhängigen Referenzimplementierung nachgeprüft. Für die Dreikörperenergie $T_\nu$ ergibt sich in 2D ein Phasendiagramm mit drei Phasen: hexagonal $\to$ Quadrat bei $\nu^*=3{,}9183649026$ (erste Ordnung) und Quadrat $\to$ Rechteck bei $\nu_2\approx 8{,}606$ (kontinuierlich). In 3D ist BCC bei allen untersuchten Exponenten das beste gefundene Gitter, während FCC ab $\nu\approx 3{,}75$ instabil wird. Planare modulare Graphfunktionen der Stringtheorie lassen sich als Graph-Zetas des dualen Graphen schreiben; zwei bekannte Identitäten wurden damit auf $10^{-14}$ bis $10^{-15}$ bestätigt.
Eine gezielte Literaturprüfung fand für den 2D-Phasenverlauf, die Minima der zugehörigen modularen Graphfunktionen $C_{s,s,s}$ und das globale 3D-Ergebnis keine Vorarbeit. Die engste Vorarbeit (arXiv:2504.07338) behandelt nur $\nu=3$ und den Bain-Pfad; ihre veröffentlichten Werte stimmen mit unserer Rechnung überein.
